% GENERATED FILE. Edit canonical lessons, then run npm run generate:books.
% canonical-source-hash: fnv1a64:a8ef399570f9ecee
% canonical-lessons: TE-C404-pedarikam, TE-C404-dhanika, TE-C404-sampada, TE-C404-svatantra, TE-C404-nirudyogam

\chapter{Rich and Poor}
\label{ch:te-rich-and-poor}
\hlchaptermodality{404}

\begin{chapteropening}
\textbf{By the end of this chapter:} \emph{I can talk about poverty, wealth and unemployment.}

Complete the last lesson of chapter 404: I can talk about poverty, wealth and unemployment.
\end{chapteropening}

\section[\texorpdfstring{\te{పేదరికం} (pēdarikaṁ)}{పేదరికం (pēdarikaṁ)}]{\te{పేదరికం} (pēdarikaṁ) — poverty}

\label{lesson:TE-C404-pedarikam}



\begin{quote}
Before the new one: say the Telugu for a scheme, then the Telugu for welfare.
\end{quote}

\subsection*{You'll want to know: \te{పేదరికం}}
\textbf{\te{పేదరికం}} — \emph{pēdarikaṁ} — \textquotedblleft{}poverty\textquotedblright{}.

\begin{quote}
\te{పేదరికం} \te{వల్ల} \te{చాలా} \te{మంది} \te{పిల్లలు} \te{చదువుకోలేరు}.
\end{quote}

\begin{quote}
\emph{pēdarikaṁ valla cālā mandi pillalu caduvukōlēru.}
\end{quote}

\begin{quote}
Because of poverty many children cannot go to school.
\end{quote}

\te{పేద}, poor, which you know, + -\te{రికం}.

\begin{cousinweb}
\textbf{\te{పేదరికం}} (\emph{pēdarikaṁ}) is native Dravidian: DEDR 4437 lists Telugu \emph{pēda}, \textquotedblleft{}poor\textquotedblright{}, and \emph{pēdaṟikamu}, \textquotedblleft{}poverty\textquotedblright{}.
\end{cousinweb}

\begin{grammarlens}[title={What you've built}]
One more word for giving reasons and opinions.
\end{grammarlens}

\subsection*{Your turn}
\begin{itemize}
\raggedright
  \item \emph{Say it:} \emph{pēdarikaṁ}
  \item \emph{Say it:} \emph{pēdarikaṁ}, once more
  \item \emph{Say it:} \emph{pēdarikaṁ valla cālā mandi pillalu caduvukōlēru.} — Because of poverty many children cannot go to school.
  \item \emph{Recall:} say the Telugu for a scheme, then the Telugu for welfare, then say \emph{pēdarikaṁ} again
\end{itemize}

\begin{tcolorbox}[breakable,colback=teal!4,colframe=teal!35!black,title={Before you move on}]
What does \te{పేదరికం} mean? (\textquotedblleft{}Poverty\textquotedblright{}.) Say it once more.
\end{tcolorbox}
\section[\texorpdfstring{\te{ధనిక} (dhanika)}{ధనిక (dhanika)}]{\te{ధనిక} (dhanika) — rich; wealthy}

\label{lesson:TE-C404-dhanika}



\begin{quote}
Before the new one: say the Telugu for welfare, then the Telugu for poverty.
\end{quote}

\subsection*{You'll want to know: \te{ధనిక}}
\textbf{\te{ధనిక}} — \emph{dhanika} — \textquotedblleft{}rich; wealthy\textquotedblright{}.

\begin{quote}
\te{ధనిక} \te{దేశాలు} \te{పేద} \te{దేశాలకు} \te{సహాయం} \te{చేయాలి}.
\end{quote}

\begin{quote}
\emph{dhanika dēśālu pēda dēśālaku sahāyaṁ cēyāli.}
\end{quote}

\begin{quote}
Rich countries should help poor countries.
\end{quote}

Before a noun: \te{ధనిక} \te{కుటుంబం}, a rich family.

\begin{cousinweb}
\textbf{\te{ధనిక}} (\emph{dhanika}) is borrowed from Sanskrit \emph{dhanika}, \textquotedblleft{}wealthy, opulent\textquotedblright{} (Monier-Williams), from \emph{dhana}, \textquotedblleft{}wealth\textquotedblright{}.
\end{cousinweb}

\begin{grammarlens}[title={What you've built}]
One more word for giving reasons and opinions.
\end{grammarlens}

\subsection*{Your turn}
\begin{itemize}
\raggedright
  \item \emph{Say it:} \emph{dhanika}
  \item \emph{Say it:} \emph{dhanika}, once more
  \item \emph{Say it:} \emph{dhanika dēśālu pēda dēśālaku sahāyaṁ cēyāli.} — Rich countries should help poor countries.
  \item \emph{Recall:} say the Telugu for welfare, then the Telugu for poverty, then say \emph{dhanika} again
\end{itemize}

\begin{tcolorbox}[breakable,colback=teal!4,colframe=teal!35!black,title={Before you move on}]
What does \te{ధనిక} mean? (\textquotedblleft{}Rich; wealthy\textquotedblright{}.) Say it once more.
\end{tcolorbox}
\section[\texorpdfstring{\te{సంపద} (sampada)}{సంపద (sampada)}]{\te{సంపద} (sampada) — wealth; riches}

\label{lesson:TE-C404-sampada}



\begin{quote}
Before the new one: say the Telugu for poverty, then the Telugu for rich.
\end{quote}

\subsection*{You'll want to know: \te{సంపద}}
\textbf{\te{సంపద}} — \emph{sampada} — \textquotedblleft{}wealth; riches\textquotedblright{}.

\begin{quote}
\te{ఆ} \te{కుటుంబం} \te{దగ్గర} \te{చాలా} \te{సంపద} \te{ఉంది}.
\end{quote}

\begin{quote}
\emph{ā kuṭumbaṁ daggara cālā sampada undi.}
\end{quote}

\begin{quote}
That family has a great deal of wealth.
\end{quote}

Also a country’s riches: \te{సహజ} \te{సంపద}, natural resources.

\begin{cousinweb}
\textbf{\te{సంపద}} (\emph{sampada}) is borrowed from Sanskrit \emph{sampad}, \textquotedblleft{}success\textquotedblright{} (CDIAL 12932), which Monier-Williams also glosses \textquotedblleft{}accomplishment, perfection\textquotedblright{}; Telugu uses it for wealth.
\end{cousinweb}

\begin{grammarlens}[title={What you've built}]
One more word for giving reasons and opinions.
\end{grammarlens}

\subsection*{Your turn}
\begin{itemize}
\raggedright
  \item \emph{Say it:} \emph{sampada}
  \item \emph{Say it:} \emph{sampada}, once more
  \item \emph{Say it:} \emph{ā kuṭumbaṁ daggara cālā sampada undi.} — That family has a great deal of wealth.
  \item \emph{Recall:} say the Telugu for poverty, then the Telugu for rich, then say \emph{sampada} again
\end{itemize}

\begin{tcolorbox}[breakable,colback=teal!4,colframe=teal!35!black,title={Before you move on}]
What does \te{సంపద} mean? (\textquotedblleft{}Wealth; riches\textquotedblright{}.) Say it once more.
\end{tcolorbox}
\section[\texorpdfstring{\te{స్వతంత్ర} (svatantra)}{స్వతంత్ర (svatantra)}]{\te{స్వతంత్ర} (svatantra) — independent}

\label{lesson:TE-C404-svatantra}



\begin{quote}
Before the new one: say the Telugu for rich, then the Telugu for wealth.
\end{quote}

\subsection*{You'll want to know: \te{స్వతంత్ర}}
\textbf{\te{స్వతంత్ర}} — \emph{svatantra} — \textquotedblleft{}independent\textquotedblright{}.

\begin{quote}
\te{ఉద్యోగం} \te{వచ్చాక} \te{ఆమె} \te{స్వతంత్ర} \te{జీవితం} \te{గడుపుతోంది}.
\end{quote}

\begin{quote}
\emph{udyōgaṁ vaccāka āme svatantra jīvitaṁ gaḍuputōndi.}
\end{quote}

\begin{quote}
Since getting a job she has been living an independent life.
\end{quote}

Before a noun: \te{స్వతంత్ర} \te{జీవితం}, an independent life. \te{స్వతంత్రంగా}, independently.

\begin{cousinweb}
\textbf{\te{స్వతంత్ర}} (\emph{svatantra}) is borrowed from Sanskrit \emph{svatantra}, \textquotedblleft{}self-dependent, independent, free\textquotedblright{} (Monier-Williams), from \emph{sva}, \textquotedblleft{}own\textquotedblright{}, and \emph{tantra}, \textquotedblleft{}a system, a rule\textquotedblright{}.
\end{cousinweb}

\begin{grammarlens}[title={What you've built}]
One more word for giving reasons and opinions.
\end{grammarlens}

\subsection*{Your turn}
\begin{itemize}
\raggedright
  \item \emph{Say it:} \emph{svatantra}
  \item \emph{Say it:} \emph{svatantra}, once more
  \item \emph{Say it:} \emph{udyōgaṁ vaccāka āme svatantra jīvitaṁ gaḍuputōndi.} — Since getting a job she has been living an independent life.
  \item \emph{Recall:} say the Telugu for rich, then the Telugu for wealth, then say \emph{svatantra} again
\end{itemize}

\begin{tcolorbox}[breakable,colback=teal!4,colframe=teal!35!black,title={Before you move on}]
What does \te{స్వతంత్ర} mean? (\textquotedblleft{}Independent\textquotedblright{}.) Say it once more.
\end{tcolorbox}
\section[\texorpdfstring{\te{నిరుద్యోగం} (nirudyōgaṁ)}{నిరుద్యోగం (nirudyōgaṁ)}]{\te{నిరుద్యోగం} (nirudyōgaṁ) — unemployment}

\label{lesson:TE-C404-nirudyogam}



\begin{quote}
Before the new one: say the Telugu for wealth, then the Telugu for independent.
\end{quote}

\subsection*{You'll want to know: \te{నిరుద్యోగం}}
\textbf{\te{నిరుద్యోగం}} — \emph{nirudyōgaṁ} — \textquotedblleft{}unemployment\textquotedblright{}.

\begin{quote}
\te{మా} \te{ప్రాంతంలో} \te{నిరుద్యోగం} \te{పెరిగింది}.
\end{quote}

\begin{quote}
\emph{mā prāntaṁlō nirudyōgaṁ perigindi.}
\end{quote}

\begin{quote}
Unemployment has risen in our area.
\end{quote}

\te{నిర్}-, without, + \te{ఉద్యోగం}, a job, which you know.

\begin{cousinweb}
\textbf{\te{నిరుద్యోగం}} (\emph{nirudyōgaṁ}) is built from Sanskrit \emph{nir-}, \textquotedblleft{}without\textquotedblright{}, and \emph{udyoga}, \textquotedblleft{}exertion, undertaking\textquotedblright{} (Monier-Williams), the word Telugu has as \te{ఉద్యోగం}, a job.
\end{cousinweb}

\begin{grammarlens}[title={What you've built}]
One more word for giving reasons and opinions.
\end{grammarlens}

\subsection*{Your turn}
\begin{itemize}
\raggedright
  \item \emph{Say it:} \emph{nirudyōgaṁ}
  \item \emph{Say it:} \emph{nirudyōgaṁ}, once more
  \item \emph{Say it:} \emph{mā prāntaṁlō nirudyōgaṁ perigindi.} — Unemployment has risen in our area.
  \item \emph{Recall:} say the Telugu for wealth, then the Telugu for independent, then say \emph{nirudyōgaṁ} again
\end{itemize}

\begin{tcolorbox}[breakable,colback=teal!4,colframe=teal!35!black,title={Before you move on}]
What does \te{నిరుద్యోగం} mean? (\textquotedblleft{}Unemployment\textquotedblright{}.) Say it once more.
\end{tcolorbox}
